\section{Results}
\label{sec:results}
The precision simulation of Appendix~\ref{app:precision} found no draw count that meets the prespecified criterion. With 60 fresh draws per split, the largest candidate, a two-point primary benefit is detected in 59 percent of simulated studies and practical equivalence is established in 19 to 21 percent, against the required 80 percent in both scenarios. As Section~\ref{sec:evaluation} specifies, the experiment therefore stops with a precision limitation. No classifier was trained, so the four primary contrasts, their interaction, and the breadth contrast $\Delta_G$ remain unmeasured. A supplementary manipulation check with a fixed draw count (Section~\ref{sec:check-results}) shows that the matching tolerances cannot be met either.

\subsection{Simulated success by draw count}
\label{sec:precision-results}
Each simulated primary contrast inherits the training-draw and test-patient variability of the selection contrast between five selected and five random patients \citep{queisler2026shortageSelection}. The breadth contrast inherits the variability of ten random patients with sixteen patches each minus five selected patients. Both templates combine 15 stored split--draw fits with 2,000 test-patient replicates. The simulation draws 2,000 studies per condition, so success rates carry Monte Carlo standard errors of at most 1.1 percentage points.

Success grows slowly with the number of draws and stays below 80 percent everywhere (Table~\ref{tab:precision}). Tripling the draws from 20 to 60 raises the rate of detecting a two-point primary benefit from 48--49 to 59 percent and the rate of establishing primary equivalence from at most 1 to about 20 percent. The breadth contrast, which uses an unadjusted interval, fares better but also falls short, at 74 and 48 percent. Increasing training-draw dispersion by 50 percent lowers every rate further; primary equivalence then reaches at most 4 percent.

\begin{table}[htbp]
\centering
\small
\renewcommand{\arraystretch}{1.15}
\begin{tabular}{@{}rlrrrrrr@{}}
\toprule
 & & \multicolumn{2}{c}{Primary contrasts, \%} & \multicolumn{2}{c}{Breadth contrast, \%} & \multicolumn{2}{c}{Half-width, points} \\
\cmidrule(lr){3-4}\cmidrule(lr){5-6}\cmidrule(l){7-8}
Draws & Dispersion & Benefit & Equivalence & Benefit & Equivalence & Primary & Breadth \\
\midrule
20 & Stored & 48--49 & 0--1 & 67 & 36 & 1.02 & 0.81 \\
 & $\times 1.5$ & 32--33 & 0 & 60 & 17 & 1.20 & 0.90 \\
\addlinespace
40 & Stored & 54--56 & 13--14 & 73 & 44 & 0.93 & 0.77 \\
 & $\times 1.5$ & 43--45 & 0 & 68 & 35 & 1.05 & 0.82 \\
\addlinespace
60 & Stored & 59 & 19--21 & 74 & 48 & 0.89 & 0.75 \\
 & $\times 1.5$ & 50--51 & 4 & 69 & 37 & 0.98 & 0.79 \\
\bottomrule
\end{tabular}
\caption{No draw count reaches 80 percent simulated success, the prespecified requirement for every contrast in both scenarios. Draws are fresh training draws per split. Benefit is the percentage of simulated studies whose lower interval bound exceeds one point under a true two-point contrast; equivalence is the percentage whose interval lies within $[-1,1]$ under a true zero contrast. Primary columns give the range over the four primary contrasts, which use simultaneous 95\% intervals; the breadth contrast uses a separate 95\% interval. Dispersion is the stored training-draw variability or that variability increased by 50 percent. Half-widths are medians over 300 simulated studies per condition.}
\label{tab:precision}
\end{table}

\subsection{Why additional draws do not suffice}
\label{sec:precision-floor}
The limit lies in the fixed evaluation sample. With training draws held fixed, test-patient resampling alone gives the selection contrast a standard error of 0.38 points, consistent with the interval half-width of 0.70 points reported for the selection gain \citep{queisler2026shortageSelection}. At 60 draws, the simulated standard error is 0.40 points, so additional draws have almost exhausted what they can remove. The simultaneous critical value over the four primary contrasts is 2.23, which lies between the unadjusted 1.96 and a Bonferroni value of 2.64 because the contrasts share test patients. The resulting half-width of 0.89 points consumes most of the one-point margin.

Both decision rules then depend on how far the estimate falls from the true contrast. Across simulated studies at 60 draws, this estimation error has a standard deviation of 0.41 points. Equivalence requires the whole interval within $[-1,1]$, which a zero contrast achieves only when its estimate lies within 0.11 points of zero. Detecting a two-point benefit requires the estimate to fall short of the true value by at most 0.11 points. Under a normal approximation, these conditions hold with probabilities of about 21 and 61 percent, close to the simulated rates. Even with unlimited draws, the test-patient standard error of 0.38 points and a critical value near 2.2 would leave a half-width of about 0.85 points. Equivalence would then remain below one third and benefit detection below two thirds. Under the same approximation, 60 draws would detect a benefit with 80 percent success from a true primary contrast of about 2.2 points; establishing equivalence would require a half-width of about 0.47 points, which this evaluation sample cannot provide.

\subsection{Supplementary manipulation check}
\label{sec:check-results}
The cohort search requires no classifier fits, so the manipulation check of Section~\ref{sec:check} was also run with 20 draws per split, the smallest candidate draw count, for all 30 classes. The search reaches both manipulated differences. Averaged over classes and draws, poor minus good coverage distance is 0.039 to 0.042 at both similarity levels, against the required 0.03, and high minus low similarity is 0.070, against the required 0.05. Matching fails: in every split, none of the 600 class--draw combinations meets all tolerances (Table~\ref{tab:check-results}).

\begin{table}[htbp]
\centering
\renewcommand{\arraystretch}{1.15}
\begin{tabular}{@{}lrrrr@{}}
\toprule
 & \multicolumn{2}{c}{Coverage distance} & & \\
\cmidrule(lr){2-3}
Cohort & Training & Validation & $\omega$ & $\Neff^{\omega}$ \\
\midrule
Good coverage, low similarity & 0.345 & 0.363 & $-0.060$ & 13.11 \\
Poor coverage, low similarity & 0.385 & 0.403 & $-0.060$ & 13.11 \\
Good coverage, high similarity & 0.345 & 0.361 & 0.010 & 9.69 \\
Poor coverage, high similarity & 0.385 & 0.402 & 0.010 & 9.69 \\
\addlinespace
Ten-patient match & 0.345 & 0.371 & 0.093 & 10.97 \\
\bottomrule
\end{tabular}
\caption{The search sets coverage and similarity independently on the training pool, but the ten-patient match cannot reproduce the effective support of the good-coverage, low-similarity cohort. Values are means over the 30 classes, three splits, and 20 draws. Training coverage distance uses eligible training patients as the reference set, validation coverage distance uses validation patients, and $\Neff^{\omega}$ is the class-specific similarity-adjusted effective support of Equation~\eqref{eq:support}.}
\label{tab:check-results}
\end{table}

\noindent Two tolerances fail. First, cohorts matched on training coverage drift apart on validation patients. The two cohorts at one coverage level differ by more than 0.005 in validation coverage distance in 78 to 85 percent of class--draw combinations, and the ten-patient match misses the validation coverage of the good-coverage, low-similarity cohort in 84 to 86 percent. On average the cohorts at one coverage level differ by at most 0.002 in validation coverage distance, so the drift reflects individual draws rather than a systematic offset.

Second, the ten-patient match comes within 5 percent of the effective support of the good-coverage, low-similarity cohort in no combination, reaching 10.97 against 13.11. The good-coverage level lies at a training coverage distance of 0.345, which is poorer than the 0.302 of ten random patients. Ten patients reproduce this coverage only when their deviations from the class mean align, which raises their similarity to 0.093 and lowers their effective support.

\section{Discussion and conclusion}
\label{sec:discussion}
Whether the gain from coverage-maximizing selection arises from better coverage or from lower between-patient similarity remains open, and both preliminary assessments of this protocol fail. On the TCGA-UT evaluation sample, practical equivalence at the one-point margin, which the interpretation of Section~\ref{sec:interpretation} needs to call either property negligible, is out of reach at any number of training draws. Coverage and similarity can be set independently on the training pool, but cohorts cannot be matched draw by draw on validation patients, and ten patients cannot match five on coverage and effective support at the same time.

These results constrain a follow-up in three ways. Every contrast evaluated on these test patients carries a standard error of at least 0.38 points, so questions should be framed as estimates with intervals and bounds rather than as equivalence. Coverage and similarity should enter the analysis as measured quantities rather than as matching constraints, because their averages transfer to validation patients while individual draws do not. Finally, coverage is tied to patient count beyond random sampling: selection cannot make five patients cover held-out patients as well as twenty, nor twenty as poorly as five.

\subsection{Outline of a follow-up experiment}
\label{sec:follow-up}
A follow-up experiment would estimate how much of the accuracy difference between five and ten random patients, 62.25 against 67.28 percent \citep{queisler2026effectiveSupport}, and of the 5.52-point selection gain \citep{queisler2026shortageSelection} coverage, similarity, and patient count each explain. Instead of matching cohorts, it would vary coverage and similarity on a three-by-three grid of target levels in five- and ten-patient cohorts, with coverage spanning the difference between random ten- and five-patient cohorts. Each class would receive a random grid cell within each fit, and class recall would be regressed on the achieved values with class and fit fixed effects. Conclusions would rest on intervals and upper bounds rather than on equivalence.

A search without classifier fits shows that this grid is attainable: 94 to 96 percent of cohorts reach their targets, 90 to 91 percent of grid cells are reached at both patient counts, and coverage and similarity are nearly uncorrelated. A grid extending to twenty patients fails for the reason given above, so the design leaves out the 4.48-point step from ten to twenty patients. Its main open risk is precision, because class recall is noisier than cohort accuracy. A separate protocol would assess precision and fix all criteria before any classifier is trained.
